Đề tài này hướng đến xây dựng một khung phương pháp luận (methodological framework) toàn diện, lấy kiến trúc Transformer làm nền tảng, nhằm giải quyết việc phân tích sắc thái cảm xúc của người dùng Việt Nam trên YouTube, đồng thời xử lý hiệu quả các trở ngại đặc thù của dữ liệu mạng xã hội.

Toàn bộ quá trình nghiên cứu được hệ thống hóa thành một quy trình chuẩn hóa, cô đọng trong năm yếu tố cốt lõi dưới đây:

\begin{itemize}[label=-]
    \item Bước khởi đầu tập trung trích xuất bình luận cùng siêu dữ liệu (metadata) từ danh sách video thuộc nhiều danh mục (categories) khác nhau do YouTube phân loại, vốn đã được chuẩn bị và tổng hợp từ trước, với mục tiêu thiết lập tập dữ liệu đa dạng và tính đại diện cao, đáp ứng yêu cầu phân tích đa miền (multi-domain).
    \item Dữ liệu thô sau đó trải qua chuỗi thao tác tiền xử lý chặt chẽ: lọc thô, chuẩn hóa, loại bỏ trùng lặp và lọc sâu, được thiết kế nhằm khử tối đa các nhiễu đặc thù của ngôn ngữ mạng xã hội (từ lặp, sự pha trộn ngôn ngữ,...), giúp dữ liệu trở nên "sạch" trong khi vẫn giữ nguyên cấu trúc ngữ nghĩa và manh mối cảm xúc gốc của bình luận.
    \item Gán nhãn bán tự động sử dụng mô hình \gls{llm}, kết hợp ba vòng hậu kiểm nhằm đưa ra phương án gán nhãn tự động cho dữ liệu quy mô lớn, giúp tiết kiệm thời gian và nguồn lực so với gán nhãn truyền thống.
    \item Vận dụng phương pháp Học chuyển giao (Transfer Learning) cùng xử lý chênh lệch phân phối dữ liệu để tinh chỉnh (fine-tune) các kiến trúc ngôn ngữ hiện đại thuộc họ BERT, giúp mô hình vượt qua rào cản mất cân bằng phân lớp (class imbalance), gia tăng độ thích ứng trước sự biến đổi liên tục của ngôn ngữ trực tuyến và duy trì khả năng khái quát hóa (generalization) cao trong môi trường đa miền.
    \item Chất lượng các mô hình được đánh giá bằng những chỉ số định lượng dựa trên các thang đo tiêu chuẩn. Việc đối chiếu toàn diện giữa kết quả dự đoán và biểu đồ huấn luyện không chỉ khắc họa bức tranh tổng thể về năng lực của từng kiến trúc mạng, mà còn là bằng chứng thực nghiệm vững chắc cho thấy khung phương pháp luận này đáp ứng đầy đủ trọng điểm đã trình bày tại phần \ref{sec:methodological_framework}
\end{itemize}
